% GENERATED FILE. Edit canonical lessons, then run npm run generate:books.
% canonical-source-hash: fnv1a64:594ee66f38bb4efd
% canonical-lessons: GE-C511-freibad, GE-C511-kraulen, GE-C511-tauchen, GE-C511-paddeln, GE-C511-kanu

\chapter{Im Wasser}
\label{ch:ge-im-wasser}
\hlchaptermodality{511}

\begin{chapteropening}
\textbf{By the end of this chapter:} \emph{I can talk about water sports.}

Complete the last lesson of chapter 511: I can talk about water sports.
\end{chapteropening}

\section[das Freibad]{das Freibad — open-air pool}

\label{lesson:GE-C511-freibad}



\begin{quote}
Before the new one: say the German for cheering, then the German for to boo.
\end{quote}

\subsection*{You'll want to know: das Freibad}
\textbf{das Freibad} — \textquotedblleft{}open-air pool\textquotedblright{}.

\begin{quote}
Im Sommer gehen wir jeden Tag ins Freibad.
\end{quote}

\begin{quote}
\emph{In summer we go to the open-air pool every day.}
\end{quote}

Plural: \emph{die Freibäder}. An indoor pool is \emph{das Hallenbad}.

\begin{cousinweb}
\emph{Freibad} is built on \emph{frei} \textquotedblleft{}free, open\textquotedblright{} and \emph{Bad} \textquotedblleft{}bath\textquotedblright{}, inherited from Old High German \emph{bad}, from Proto-Germanic \emph{baþą} (EtymDB).
\end{cousinweb}

\begin{grammarlens}[title={What you've built}]
One more word for telling, explaining and describing at B1.
\end{grammarlens}

\subsection*{Your turn}
\begin{itemize}
\raggedright
  \item \emph{Say it:} \emph{das Freibad}
  \item \emph{Say it:} \emph{das Freibad}, once more
  \item \emph{Say it:} a sentence of your own with \emph{das Freibad}
  \item \emph{Recall:} say the German for cheering, then the German for to boo, then say \emph{das Freibad} again
\end{itemize}

\begin{tcolorbox}[breakable,colback=teal!4,colframe=teal!35!black,title={Before you move on}]
What does das Freibad mean? (\textquotedblleft{}Open-air pool\textquotedblright{}.) Say it once more.
\end{tcolorbox}
\section[kraulen]{kraulen — to swim front crawl}

\label{lesson:GE-C511-kraulen}



\begin{quote}
Before the new one: say the German for to boo, then the German for open-air pool.
\end{quote}

\subsection*{You'll want to know: kraulen}
\textbf{kraulen} — \textquotedblleft{}to swim front crawl\textquotedblright{}.

\begin{quote}
Mein Sohn kann schon kraulen.
\end{quote}

\begin{quote}
\emph{My son can already swim front crawl.}
\end{quote}

A different verb \emph{kraulen} means to scratch gently: \emph{den Hund hinter den Ohren kraulen}.

\begin{cousinweb}
The swimming verb \emph{kraulen} was borrowed from English \emph{crawl}; the other \emph{kraulen}, \textquotedblleft{}to scratch gently\textquotedblright{}, comes from \emph{krauen} (EtymDB).
\end{cousinweb}

\begin{grammarlens}[title={What you've built}]
One more word for telling, explaining and describing at B1.
\end{grammarlens}

\subsection*{Your turn}
\begin{itemize}
\raggedright
  \item \emph{Say it:} \emph{kraulen}
  \item \emph{Say it:} \emph{kraulen}, once more
  \item \emph{Say it:} a sentence of your own with \emph{kraulen}
  \item \emph{Recall:} say the German for to boo, then the German for open-air pool, then say \emph{kraulen} again
\end{itemize}

\begin{tcolorbox}[breakable,colback=teal!4,colframe=teal!35!black,title={Before you move on}]
What does kraulen mean? (\textquotedblleft{}To swim front crawl\textquotedblright{}.) Say it once more.
\end{tcolorbox}
\section[tauchen]{tauchen — to dive (underwater)}

\label{lesson:GE-C511-tauchen}



\begin{quote}
Before the new one: say the German for open-air pool, then the German for to swim front crawl.
\end{quote}

\subsection*{You'll want to know: tauchen}
\textbf{tauchen} — \textquotedblleft{}to dive (underwater)\textquotedblright{}.

\begin{quote}
Im Urlaub sind wir zwischen bunten Fischen getaucht.
\end{quote}

\begin{quote}
\emph{On holiday we dived among colourful fish.}
\end{quote}

Perfekt mostly with \emph{sein}: \emph{wir sind getaucht}. A diver is \emph{der Taucher}.

\begin{cousinweb}
\emph{tauchen} is inherited from Middle High German \emph{tūchen}, \emph{touchen}, from Proto-Germanic \emph{dūkaną} (EtymDB).
\end{cousinweb}

\begin{grammarlens}[title={What you've built}]
One more word for telling, explaining and describing at B1.
\end{grammarlens}

\subsection*{Your turn}
\begin{itemize}
\raggedright
  \item \emph{Say it:} \emph{tauchen}
  \item \emph{Say it:} \emph{tauchen}, once more
  \item \emph{Say it:} a sentence of your own with \emph{tauchen}
  \item \emph{Recall:} say the German for open-air pool, then the German for to swim front crawl, then say \emph{tauchen} again
\end{itemize}

\begin{tcolorbox}[breakable,colback=teal!4,colframe=teal!35!black,title={Before you move on}]
What does tauchen mean? (\textquotedblleft{}To dive (underwater)\textquotedblright{}.) Say it once more.
\end{tcolorbox}
\section[paddeln]{paddeln — to paddle}

\label{lesson:GE-C511-paddeln}



\begin{quote}
Before the new one: say the German for to swim front crawl, then the German for to dive.
\end{quote}

\subsection*{You'll want to know: paddeln}
\textbf{paddeln} — \textquotedblleft{}to paddle\textquotedblright{}.

\begin{quote}
Wir paddeln mit dem Boot über den See.
\end{quote}

\begin{quote}
\emph{We paddle across the lake in the boat.}
\end{quote}

Perfekt with \emph{sein} for a trip: \emph{Wir sind über den See gepaddelt.}

\begin{cousinweb}
\emph{paddeln} was borrowed from English \emph{to paddle}; EtymDB has no entry for the German verb.
\end{cousinweb}

\begin{grammarlens}[title={What you've built}]
One more word for telling, explaining and describing at B1.
\end{grammarlens}

\subsection*{Your turn}
\begin{itemize}
\raggedright
  \item \emph{Say it:} \emph{paddeln}
  \item \emph{Say it:} \emph{paddeln}, once more
  \item \emph{Say it:} a sentence of your own with \emph{paddeln}
  \item \emph{Recall:} say the German for to swim front crawl, then the German for to dive, then say \emph{paddeln} again
\end{itemize}

\begin{tcolorbox}[breakable,colback=teal!4,colframe=teal!35!black,title={Before you move on}]
What does paddeln mean? (\textquotedblleft{}To paddle\textquotedblright{}.) Say it once more.
\end{tcolorbox}
\section[das Kanu]{das Kanu — canoe}

\label{lesson:GE-C511-kanu}



\begin{quote}
Before the new one: say the German for to dive, then the German for to paddle.
\end{quote}

\subsection*{You'll want to know: das Kanu}
\textbf{das Kanu} — \textquotedblleft{}canoe\textquotedblright{}.

\begin{quote}
Mit dem Kanu sind wir den Fluss hinuntergefahren.
\end{quote}

\begin{quote}
\emph{We went down the river by canoe.}
\end{quote}

Plural: \emph{die Kanus}.

\begin{cousinweb}
\emph{Kanu} came into German through English \emph{canoe} and Spanish \emph{canoa}, usually traced to a Caribbean word for a dugout boat; EtymDB has no entry for the German word.
\end{cousinweb}

\begin{grammarlens}[title={What you've built}]
One more word for telling, explaining and describing at B1.
\end{grammarlens}

\subsection*{Your turn}
\begin{itemize}
\raggedright
  \item \emph{Say it:} \emph{das Kanu}
  \item \emph{Say it:} \emph{das Kanu}, once more
  \item \emph{Say it:} a sentence of your own with \emph{das Kanu}
  \item \emph{Recall:} say the German for to dive, then the German for to paddle, then say \emph{das Kanu} again
\end{itemize}

\begin{tcolorbox}[breakable,colback=teal!4,colframe=teal!35!black,title={Before you move on}]
What does das Kanu mean? (\textquotedblleft{}Canoe\textquotedblright{}.) Say it once more.
\end{tcolorbox}
